\section{Discussion and Limitations}
\label{sec:discussion}

\paragraph{What VLMs seem to see.}
Ordinary-Bench is designed to test a narrower and more diagnostic claim than general multimodal competence. A model may identify objects correctly, produce plausible captions, or answer some spatial prompts accurately while still failing to maintain a stable geometric belief over the full scene. The key distinction is therefore not between ``correct'' and ``incorrect'' in the abstract, but between local success and globally coherent, visually grounded structure.

\paragraph{Why scene-level analysis matters.}
This distinction is important because many downstream uses of VLMs implicitly assume more than isolated correctness. Planning, embodied interaction, and multi-step reasoning all require judgments that are mutually compatible. If a model says that one pair of objects is closer than another, and separately answers directional probes about the same scene, those answers should support a single world rather than a patchwork of unrelated guesses. Scene belief reconstruction makes this compatibility test explicit.

\paragraph{Why controlled synthetic data is a feature here.}
The benchmark deliberately uses controlled rendered scenes instead of unconstrained real-world imagery. This lowers ecological realism, but it makes exact geometry, dense probing, and visual perturbation possible. For the scientific question posed here, this tradeoff is appropriate: it is difficult to separate visual evidence from language priors, or to measure scene-level coherence, when the underlying scene structure is itself ambiguous. Controlled synthetic data therefore acts as an instrument for diagnosis rather than as a claim of direct real-world coverage.

\paragraph{Limitations.}
The current benchmark has several limitations. First, the core reconstruction pipeline operates in 2D or 2.5D terms even though scenes are rendered from 3D environments; this matches the semantics of the current TRR implementation but does not amount to full 3D scene recovery. Second, the main evaluation setting treats models as API-level black boxes, so we analyze externalized behavior rather than hidden representations. Third, the benchmark focuses on simple geometric objects in synthetic scenes; transferring conclusions to natural images or embodied environments requires additional validation. Finally, the difficulty of fine-grained TRR prediction means that some conclusions may be clearer at coarser directional granularity than at exact clock-hour resolution.

\paragraph{Future directions.}
The natural next steps are to complete the temporal extension already prototyped in the repository, expand the release package for broader community use, and connect behavioral reconstruction with internal representation analysis for open-weight models. Another promising direction is to combine selective answering with richer output forms such as set-valued or graph-structured scene descriptions, which could yield cleaner constraint extraction than forced single-label responses.
